\documentclass{article}
\usepackage{amsfonts, amsmath, amsthm, amssymb, enumerate, fancyhdr, gensymb, lastpage, mathtools, parskip, graphicx}
\usepackage{xcolor, tikz-cd}
\usepackage[a4paper, total={6in, 8in}]{geometry}
\usepackage{setspace}

\usepackage[OT2,T1]{fontenc}
\DeclareSymbolFont{cyrletters}{OT2}{wncyr}{m}{n}
\DeclareMathSymbol{\Sha}{\mathalpha}{cyrletters}{"58}

% Theorem environments
\newtheorem{theorem}{Theorem}
\newtheorem{lemma}[theorem]{Lemma}
\newtheorem{proposition}[theorem]{Proposition}
\newtheorem{corollary}[theorem]{Corollary}

% Definition-style environments
\theoremstyle{definition}
\newtheorem{definition}[theorem]{Definition}
\newtheorem{example}[theorem]{Example}

% Remark-style environments
\theoremstyle{remark}
\newtheorem{remark}[theorem]{Remark}


%\onehalfspacing
\setstretch{1.3}
\newcommand{\wg}[1]{\textcolor{violet}{#1}}

%\renewcommand{\C}{\mathcal C}
\newcommand{\Z}{\mathbb Z}
\newcommand{\Ztwo}{\mathbb Z / 2\mathbb Z}
\newcommand{\Ztwosq}{(\mathbb Z / 2 \mathbb Z)^2}
\newcommand{\GalK}{G_{\bar K / K}}
\newcommand{\Q}{\mathbb Q}
\newcommand{\R}{\mathbb R}
\newcommand{\im}{\text{im }}
\renewcommand{\div}{\, \vert \,}
\newcommand{\F}{\mathbb F}
\newcommand{\EQW}{E'(\Q)/2E(\Q)}
\newcommand{\Etors}{E(\Q)_{\text{tors}}}
\renewcommand{\F}[1]{\mathbb F_{#1}}
\newcommand{\Qst}{\mathbb Q(S,2)}
\newcommand{\Sphi}{S^\phi}
\newcommand{\Sphihat}{S^{\hat \phi}}
\newcommand{\bark}{{\bar{k}}}
\newcommand{\Qsqs}{(\Q^\times)/(\Q^\times)^2}
\newcommand{\Zn}[1]{\mathbb Z / #1 \mathbb Z}

\newenvironment{wgenv}{\color{violet}}{}

\newcommand{\smalltwobytwo}[4]{
\left( \begin{smallmatrix*}
   #1 & #2 \\ #3 & #4 
\end{smallmatrix*}\right) 
}
\newcommand{\bmat}[1]{
    \begin{bmatrix*}
        #1
    \end{bmatrix*}
}

\title{
    Free Resolutions
}
\author{Will Gilroy}
\date{July 2026}

\begin{document}
\maketitle

This is a note discussing the basic construction of Free Resolutions and the
computation of free resolutions for various modules.

\section{Basic Construction}
Let $R$ be a commutative ring with $1$ and let $M$ be a finitely generated
module over $R$.
If $M$ is generated by $x_1, \cdots, x_n$ elements with $m$ relations $f_i$ between the generators
then recall that we can encode the presentation of $M$ with an exact sequence
\[
    R^m \to R^n \to M \to 0
\]
where $R^n \to M$ is given by $a_i \to x_i$ and where $R^m \to R^n$ is given by
$b_i \to f_i(a_i)$. 
This is already pretty good. However, we can go further. Recalling that $R$-mod
is an abelian category gives us that $\ker(R^m \to R^n)$ is also an $R$-module.
And so we can ask for a presentation of $\ker(R^n \to M)$. This gives us a
further exact sequence
\[
    R^{m_2} \to R^{n_2} \to R^m \to R^n \to M \to 0. 
\]
Then, we are able to iterate this procedure until we no longer have any relations
between the set of generators of the current iterate. This process may not
terminate \wg{(even when $M$ is fg? we should find an example)}.
Repeating this procedure gives the \textit{free resolution of $M$}
\[
    \cdots \to R^{k_n} \to \cdots R^{k_2} \to R^{k_1} \to R^{k_0} \to M \to 0.
\]
This is called 
``Free'' since every object in the chain (except the final one) is a free
module. 
Below we will compute some examples of free resolutions of various modules.
\footnote{
\wg{I am curious to note some basic thoery here. I suppose it is clear that a 
fg module should always have a free resolution. Is it clear that a given module
has a unique free resolution (I suppose the above procedure can only spit out
one exact sequence, but if we define a free resolution to be an exact sequence
whose tail is the presentation then is there only one such sequence?)}
}

\section{Examples}
These examples are sorted into an order of increasing complexity. However, 
in all the following examples, $R$ is a polynomial ring and $M$ is a quotient of
$R$. One can check these examples in Macaulay2 with the code
\wg{what's the use of verbatim}

\begin{example}
    Let $R = k[x]$ and $M = R/(x^n)$ for some fixed $n$. 
    Then notice that $M$ is generated by $\bar 1 \in R/(x^n)$ as an $R$-module. 
    Moreover, we only have a single relation $x^n \cdot \bar 1 = 0$. 
    We can then capture this in the presentation
    \[
        R^1 \xrightarrow{1 \mapsto x^n} 
        R^1 \xrightarrow{1 \mapsto \bar 1}
        M \to 0.
    \]
    The kernel of $R^1 \to M$ is exactly the ideal $(x^n)$ (as expected).
    Interpreting $(x^n)$ as an $R$-module we find that it is a rank $1$ free
    module (generated by $x^n$). That is a presentation for $(x^n)$ as an
    $R$-module is 
    \[
        0 \to R^1 \xrightarrow{1 \mapsto x^n} (x^n).
    \]
    \wg{Is the above sequence correct?}
    Now since $\ker(R^1 \to (x^n)) = 0$ we can no longer iterate our procedure.
    This leaves us with a \wg{(the?)} free resolution 
    \[
        R^1 \xrightarrow{1 \mapsto x^n}
        R^1 \xrightarrow{1 \mapsto x^n}
        R^1 \xrightarrow{1 \mapsto \bar 1}
        M \to 0. 
    \]
    The final map does not give us any additional information \wg{(what do I
    mean by this? It's clear that we should trim the sequence however)}
    and so we actually call the following sequence the free resolution of $M =
    R/(x^n)$. 
    \[
        R^1 \xrightarrow{1 \mapsto x^n}
        R^1 \xrightarrow{1 \mapsto \bar 1}
        M \to 0. 
    \]
    This example shows us that whenever our kernel module is free we are done.
\end{example}

\begin{example}
    Let $R = k[x,y]$ and $M = R/(x^2, y^2)$ elements in $M$ are of the form 
    $f = a_0 + a_1 x + a_2 y + a_3 xy$. As an $R$-module $M$ is generated by
    $\bar 1 \in M$. We have the relations $x^2 = 0$ and $y^2 = 0$. 
    \wg{We need to describe what it means for these relations to span all the kernel relations }
    A presentation for $M$ is thus
    \[
        R^2 \xrightarrow{\bmat{x^2 & y^2}}
        R^1 \xrightarrow{1 \to \bar 1}
        M \to 0
    \]
    Now, the kernel module $\ker(R^1 \to M) = (x^2, y^2)$, interpreting the
    ideal as an $R$-module.
    This module is generated by two elements $x^2, y^2$ and we have the
    $R$-linear relation that $y^2(x^2) - x^2(y^2) = 0$. 
    Thus our next iterated sequence is
    \[
        R^1 \xrightarrow{\bmat{y^2 \\ -x^2}}
        R^2 \xrightarrow{\bmat{x^2 & y^2}}
        R^1 \xrightarrow{1 \to \bar 1}
        M \to 0. 
    \]
    Moreover the kernel module $\ker(R^2 \to (x^2, y^2)) = 0$ and so we are
    done. The above sequence is the free resolution for $M = R/(x^2, y^2)$. 
\end{example}

\begin{example}
    
\end{example}

\end{document}